% !TEX root = ../main.tex

\section{Related Work}
\label{sec:related}

\noindent\textbf{Continual Learning in Task Streams.}
Continual learning studies systems that learn from a sequence of tasks while retaining earlier capabilities~\cite{buzzega2020dark,wang2022learning,wang2022dualprompt}. Domain and test-time adaptation address distribution shift~\cite{ben-david2010da,ganin2016dann,wang2021tent,liang2020shot}, and mixture-of-experts methods route heterogeneous inputs to specialised components~\cite{jacobs1991moe,shazeer2017moe}. These directions cover pieces of our D1--D3 setting, but they usually adapt model weights, classifiers, or expert modules. In contrast, we study the harness-level analogue.

\noindent\textbf{Self-improving and self-evolving agents.}
The closest LLM-agent precedents to our setting are systems that update the harness directly from execution feedback. A-Evolve~\cite{lin2026position} introduces a linear-chain evolver: each cycle reads batch trajectories and mutates prompts, skills, memory, and tools. GEPA~\cite{agrawal2025gepa} adds reflective Pareto prompt evolution using textual feedback rather than scalar reward. Meta-Harness~\cite{lee2026meta} uses a growing filesystem archive and Claude Code as the proposer. Continual Harness~\cite{karten2026continual} enables online adaptation within a single continuous deployment run through alternating action/refinement cycles. SkillOS~\cite{ouyang2026skillos} learns a skill curation policy via reinforcement learning, training a curator to select and refine reusable skills from repeated interactions. 
